\chapter*{Introducción}
\addcontentsline{toc}{chapter}{Introducción}

Las encuestas de hogares miden lo que ningún registro administrativo cubre por completo, como la condición de actividad de la población, pero no cubren todo el territorio. Cada dominio para el que se quiere publicar una cifra precisa exige un número mínimo de hogares encuestados, y ese número cuesta en encuestadores, desplazamientos y supervisión. Por eso las encuestas concentran la muestra en los territorios de mayor peso. En Colombia, la Gran Encuesta Integrada de Hogares (GEIH), que realiza el Departamento Administrativo Nacional de Estadística (DANE), es la fuente oficial del mercado laboral, y en 2018 publicaba la tasa de desempleo de 23 ciudades capitales y áreas metropolitanas (A.M.), además de los totales nacional, por zona y por departamento \cite[p.~9]{DANE_GEIH_Metodologia2016}.

El resultado es que no se conoce la tasa de desempleo de los municipios medianos y pequeños. En el Cauca y el Valle del Cauca, 81 de los 84 municipios no tienen ninguna estimación directa, y los tres que la tienen se miden dentro de una ciudad o un área metropolitana. Las políticas de empleo de esos municipios se diseñan con la cifra del departamento o de la capital, aunque la evidencia para Colombia muestra que el desempleo municipal difiere mucho de un lugar a otro y que esas diferencias persisten \cite[pp.~124--125]{Diaz2016}.

La estimación en áreas pequeñas (SAE, por sus siglas en inglés) se desarrolló para este tipo de vacío: combina la encuesta con información auxiliar disponible para todo el territorio y, con un modelo ajustado a los dominios que sí tienen muestra, estima el indicador en los que no la tienen \cite[p.~17]{Molina2019}. Este trabajo usó el modelo de Fay-Herriot \cite{FayHerriot1979}, que tiene dos partes. El modelo de muestreo trata la tasa directa de la GEIH de cada dominio como una medición con error de su valor verdadero, con una varianza conocida. El modelo de enlace relaciona ese valor verdadero con covariables municipales de TerriData, el sistema de estadísticas territoriales del Departamento Nacional de Planeación (DNP), más un efecto aleatorio de área. Se probó además una extensión espacial, en la que los efectos de áreas cercanas pueden estar correlacionados.

La metodología se organizó en 16 etapas. Se integraron la GEIH de 2018, TerriData, la división político-administrativa del DANE y las cifras publicadas por el DANE; se estimó la tasa de desempleo anual de 23 dominios (ciudades o ciudades con su área metropolitana) con su varianza por bootstrap de viviendas; se seleccionaron covariables de un catálogo construido con la literatura mediante un Fay-Herriot univariado y un control de falsos descubrimientos; se eligió un modelo Fay-Herriot clásico como referencia y se comparó con ocho especificaciones espaciales en condiciones controladas; se validó la estimación de cada dominio retirándolo del ajuste; y se estimó la tasa de 83 municipios del Cauca y del Valle del Cauca, con su error, su estado de extrapolación y un ajuste para que las estimaciones agregadas reprodujeran las cifras directas.

Las tasas directas de los 23 dominios reprodujeron las del DANE y todas resultaron confiables, con coeficientes de variación entre 2.01~\% y 3.23~\%. De 18 covariables candidatas, solo la tasa de homicidios quedó en el modelo, con una asociación positiva que depende de un dominio, Quibdó; al repetir la selección sin cada dominio, el procedimiento estimó peor que un modelo sin covariables. No se encontró evidencia de dependencia espacial: ninguna especificación espacial tuvo una prueba de razón de verosimilitudes significativa y las estimaciones de los dominios coincidieron con las del modelo clásico. Ninguno de los 83 municipios quedó con una estimación confiable, 21 de ellos son extrapolación del modelo, y el patrón territorial de las estimaciones reproduce el de la tasa de homicidios. El aporte del trabajo es, por tanto, la metodología: un procedimiento reproducible que informa con qué calidad y bajo qué supuestos se puede estimar el desempleo de un municipio sin encuesta, y que puede repetirse con mejores covariables.

El documento se organiza así. La Sección~\ref{sec:problema} plantea y formula el problema y la Sección~\ref{sec:objetivos} presenta los objetivos. La Sección~\ref{sec:marco-referencia} reúne el marco teórico y los antecedentes. La Sección~\ref{sec:datos} describe las fuentes de información y la Sección~\ref{sec:metodologia}, las 16 etapas de la metodología. La Sección~\ref{sec:resultados} presenta los resultados de cada etapa en el mismo orden. La Sección~\ref{sec:discusion} discute su alcance y el trabajo futuro, y la Sección~\ref{sec:conclusion} cierra con las conclusiones por objetivo.
